\chapter{Quantum computation}
The main idea for using Majorana fermions for quantum computation is through a process called braiding. In this process the Majoranas are spatially rotated around eachother, producing a difference in phase. In 1D this is however hard to do; rotating the whole chain around itself has no effect. And if two Majoranas meet, they simply combine into a normal fermion, allowing the stored information to decohere. Several methods are being discussed to overcome this obstacle, and we will study two of them. 

\section{Logical qubits from Majorana modes}

\subsection{Degenerate ground states}
Recalling the previously discussed Kitaev chain Hamiltonian, and the following definition using the Majorana operators $c_{m}$ with properties 
\begin{align}
    d_j &= \frac{1}{2} \left(c_{2j} + i c_{2j+1} \right) \\
    d_j^\dagger &= \frac{1}{2} \left(c_{2j} - i c_{2j+1} \right)\\
    c_{m}^\dagger &= c_m, \hspace{4mm} \{c_l, c_m\} = 2\delta_{lm}.
\end{align}
We recall that in the trivial phase, Majoranas on all sites pair up into ordinary fermions with non-zero energy. However in the topological phase the chain experiences unpaired Majoranas at the edges, which we want well-separated (meaning their wavefunctions overlap negligibly). For a 1D wire of $2M$ such edge Majoranas one can construct $M$ fermion operators as
\begin{align}
    f_j &= \frac{1}{2} \left(c_{2j} + i c_{2j+1} \right).
\end{align}
These correspond to zero-energy modes, and can give rise to $2^M$ degenerate ground states which can be labeled by occupation numbers $n_j = f_j^\dagger f_j = 0, 1$. 

Wave function representation

Conservation of parity

Four Majoranas - a proper qubit

Nonlocality - spatial separation of Majoranas

\cite{aliceaNonAbelianStatisticsTopological2011a} \cite{aliceaNewDirectionsPursuit2012}

\section{Braiding and non-abelian statistics (exchange of Majoranas)}

\subsection{General result}
Berry phase

\subsection{Clifford group}
?

\subsection{What's missing}
Something something non-topological 

supplementary operations

\section{Methods using the Kitaev chain}


\subsection{T-junction}
The T-junction is a proposed architecture for allowing braiding with a 1D Kitaev chain. Simply put, the idea stems from connecting a 1D chain to the middle of another, producing a T-shape. In this geometry, one can in theory braid the Majoranas by altering the parameters of the Hamiltonian to rotate the edge states around eachother without touching. 

\begin{figure}
    \centering
    % \includegraphics{T-junction}
    \caption{}
    \label{fig:T_junction}
\end{figure}

\subsubsection{Exchange of two Majoranas}

\subsubsection{Network of wires}

\subsection{Quantum dots}

\subsection{Others}
